\section{Valores, variables y expresiones}

\subsection{Valores}

Los textos en español se componen de palabras; los programas en Python se componen de \emph{valores}. Algunos ejemplos de valores son números y secuencias de caracteres, como estos:
\begin{minted}{python}
42
3.14
"Hola"
\end{minted}
El primer valor es el número entero 42; el segundo, el número decimal 3.14; y el tercero, una secuencia de cuatro caracteres del teclado: \mintinline{python}|H|, \mintinline{python}|o|, \mintinline{python}|l|, \mintinline{python}|a|.

Análogo a como las palabras tienen categorías: sustantivos, verbos, etc.; los valores tienen \emph{tipos de dato}\footnote{Hay más tipos de dato en Python, que se exponen más adelante.}:

\begin{table}[h]
\centering
\footnotesize
\begin{tabular}{|l|l|l|p{3.5cm}|l|}
\hline
\textbf{Tipo de dato} & \textbf{Nombre (inglés)} & \textbf{Traducción} & \textbf{Descripción} & \textbf{E.g.} \\
\hline
\mintinline{python}|int| & \textit{integer} & número entero & Números enteros & \mintinline{python}|42|, \mintinline{python}|-17|, \mintinline{python}|0|, \mintinline{python}|1_000_000| \\
\hline
\mintinline{python}|float| & \textit{floating point number} & número de punto flotante & Números decimales & \mintinline{python}|3.14|, \mintinline{python}|-0.5|, \mintinline{python}|.35| \\
\hline
\mintinline{python}|str| & \textit{string} & cadena de caracteres\footnotemark & Secuencias de caracteres del teclado & \mintinline{python}|"Hol4"|, \mintinline{python}|'u78%6#a'| \\
\hline
\end{tabular}
\normalsize
\caption{Algunos tipos de dato primitivos en Python}
\label{tab:tipos-dato}
\end{table}
\footnotetext{La palabra \textit{string} generalmente traduce a hilo y en algunos contextos a cadena. En el contexto de la programación, es convencional traducirla como cadena de caracteres.}

Las cadenas de caracteres deben escribirse entre comillas para que la máquina entienda dónde comienza y dónde termina la secuencia de caracteres que se está representando. Hay cuatro formas de escribirlas:

\begin{itemize}
  \item Usando comillas dobles: \mintinline{python}|"Hola"|
  \item Usando comillas simples: \mintinline{python}|'H0l4'|
  \item Usando triples comillas dobles: \mintinline{python}|"""Hola"""|
  \item Usando triples comillas simples: \mintinline{python}|'''Python'''|
\end{itemize}

Si entre los caracteres de una cadena se quiere incluir algún tipo de comillas, es sensato delimitar la cadena con otro tipo de comillas, para que no haya ambigüedad. E.g.: \mintinline{python}|"Don't!"|, representa la secuencia de caracteres de teclado \mintinline{python}|D|, \mintinline{python}|o|, \mintinline{python}|n|, \mintinline{python}|'|, \mintinline{python}|t|, \mintinline{python}|!|. Si se hubiese escrito \mintinline{python}|'Don't!'|, no es claro cuál comilla termina la cadena.

Las triples comillas (dobles o simples) permiten representar cadenas de caracteres que contienen saltos de línea. E.g.:
\begin{minted}{python}
"""Hola
 Mundo"""
\end{minted}

\subsection{Variables}

Una \emph{variable} es un nombre que se asigna a un valor. Las variables se utilizan para almacenar valores en la memoria del programa.

Para definir una variable se utiliza el \emph{operador de asignación}, que es el caracter \mintinline{python}|=|. 

\begin{tip}
  El símbolo de igualdad en matemáticas y el operador de asignación en Python utilizan el mismo caracter: \mintinline{python}|=|. Sin embargo, en cada contexto significa cosas diferentes.\\
  
  En programación sucede bastante que se utilizan caracteres que en algún otro contexto significan otra cosa.
\end{tip}

La sintaxis, de izquierda a derecha, se compone de: el nombre de la variable, el operador de asignación y el valor que se le asigna a la variable. E.g.:

\begin{minted}{python}
  persona = "David"
  altura = 1.77
\end{minted}

Allí:
\begin{itemize}
  \item Se toma la secuencia de caracteres \mintinline{python}|D|, \mintinline{python}|a|, \mintinline{python}|v|, \mintinline{python}|i|, \mintinline{python}|d| y se almacena en la memoria del programa bajo el identificador \mintinline{python}|persona|.
  \item Se toma el número decimal 1.77 y se almacena en la variable \mintinline{python}|altura|.
\end{itemize}

Tiene sentido otorgar nombres a los valores para poder usarlos después. Los identificadores de las variables deben seguir las siguientes reglas:
\begin{itemize}
  \item Deben ser alfanuméricos: letras, números y guiones bajos.
  \item No pueden contener espacios; para usar varias palabras dentro del identificador, se separan con guiones bajos: \mintinline{python}|nombre_de_la_variable|.
  \item No pueden empezar por dígitos ni deben empezar por mayúsculas.
  \item No pueden ser \emph{palabras reservadas} del lenguaje, que son palabras específicas que ya tienen un significado en Python.\footnote{Usualmente en este punto los libros de texto listan las palabras reservadas. Yo no lo hago porque, primero, no es necesario memorizarlas; y, segundo, es bastante improbable que el nombre de una variable preciso coincida con una palabra reservada.}
\end{itemize}

\subsubsection{Sobreescritura y tipado dinámico}

Se puede sobreescribir el valor de una variable por otro y la máquina olvida el valor anterior. E.g.:

\begin{minted}{python}
  persona = 3
\end{minted}

Ahora, la variable \mintinline{python}|persona| tiene el valor \mintinline{python}|3| y deja de hacer referencia el valor anterior, \mintinline{python}|"David"|.

Además, cambió el tipo de dato de la variable \mintinline{python}|persona|, que antes era \mintinline{python}|str| y ahora es \mintinline{python}|int|. Cambiar el tipo de dato de una variable es posible en Python, pero no lo es en muchos otros lenguajes de programación. A esa característica de Python se le llama \emph{tipado dinámico}\footnote{Lo descrito en este párrafo es realmente una consecuencia del tipado dinámico, no una definición. Una definición escapa el alcance de estos apuntes, pero podría ser: un lenguaje es dinámicamente tipado si su sistema de tipos realiza verificaciones en tiempo de ejecución, en lugar de en tiempo de compilación.}.

\subsection{Expresiones}

Una \emph{expresión} es una combinación de valores y operadores que el intérprete de Python puede evaluar para obtener un valor. En matemática, se puede evaluar la expresión $1 + 1$ para obtener el valor $2$; en Python, es similar.

Los \emph{operadores} son símbolos que representan operaciones entre valores de un tipo de dato determinado. 

\subsubsection{Operadores de cadenas de caracteres}

Las operaciones que se pueden realizar con valores \mintinline{python}|str| son la concatenación y la repetición:

\begin{table}[h]
\centering
\begin{tabular}{|l|c|p{5cm}|c|c|}
\hline
\textbf{Operador} & \textbf{Símbolo} & \textbf{Descripción} & \textbf{E.g.} & \textbf{Resultado} \\
\hline
Concatenación & \mintinline{python}|+| & Une dos cadenas de caracteres & \mintinline{python}|"Hola" + "Mundo"| & \mintinline{python}|"HolaMundo"| \\
\hline
Repetición & \mintinline{python}|*| & Opera un \mintinline{python}|str| con un \mintinline{python}|int|; repite el \mintinline{python}|str| tantas veces como indica el \mintinline{python}|int| & \mintinline{python}|"Hola" * 3| & \mintinline{python}|"HolaHolaHola"| \\
\hline
\end{tabular}
\caption{Operadores de cadenas de caracteres en Python}
\label{tab:operadores-strings}
\end{table}

\subsubsection{Operadores aritméticos}

Las operaciones que se pueden realizar con valores de tipo de dato numérico, i.e. \mintinline{python}|int| o \mintinline{python}|float|, son las operaciones aritméticas convencionales:

\begin{table}[h]
\centering
\begin{tabular}{|l|c|c|c|c|}
\hline
\textbf{Operador} & \textbf{Símbolo} & \textbf{Ejemplo matemático} & \textbf{Ejemplo en Python} & \textbf{Resultado} \\
\hline
Exponenciación & \mintinline{python}|**| & $3^4$ & \mintinline{python}|3**4| & \mintinline{python}|81| \\
\hline
Multiplicación & \mintinline{python}|*| & $7 \times 6$ & \mintinline{python}|7 * 6| & \mintinline{python}|42| \\
\hline
División & \mintinline{python}|/| & $11 / 4$ & \mintinline{python}|11 / 4| & \mintinline{python}|2.75| \\
\hline
División entera & \mintinline{python}|//| & $\left\lfloor 11 / 4 \right\rfloor$ & \mintinline{python}|11 // 4| & \mintinline{python}|2| \\
\hline
Módulo & \mintinline{python}|%| & $11 \mod 4$ & \mintinline{python}|11 % 4| & \mintinline{python}|3| \\
\hline
Suma & \mintinline{python}|+| & $10 + 7$ & \mintinline{python}|10 + 7| & \mintinline{python}|17| \\
\hline
Resta & \mintinline{python}|-| & $20 - 8$ & \mintinline{python}|20 - 8| & \mintinline{python}|12| \\
\hline
\end{tabular}
\caption{Operadores aritméticos en Python}
\label{tab:operadores-aritmeticos}
\end{table}

Las operaciones se pueden combinar para construir expresiones más complejas. E.g.: \mintinline{python}|2 + 4 * 2| evalúa a \mintinline{python}|10|. Python no evalúa ingenuamente de izquierda a derecha, sino que respeta la \emph{jerarquía de operaciones}%
%
\footnote{Estos apuntes presuponen familiaridad con la jerarquía de operaciones; sin embargo, en esencia, el orden es: 1) exponenciación; 2) multiplicación y división (incluyendo división entera y módulo); y 3) suma y resta. Las operaciones con igual prioridad se evalúan de izquierda a derecha. Si existen paréntesis (o corchetes), las operaciones dentro de ellos tienen prioridad sobre las externas, y los paréntesis más internos se evalúan antes que los más externos.} %
%
clásica de la aritmética.

\paragraph{División entera y módulo.} Seguramente el único par de operadores no familiares en la tabla \ref{tab:operadores-aritmeticos} son los operadores de división entera (\mintinline{python}|//|) y módulo (\mintinline{python}|%|).

Cuando se realiza una división, las partes son:
\begin{gather*}
  \text{dividendo} \operatorname{\texttt{/}} \text{divisor} = \text{cociente}
\end{gather*}
Usualmente el cociente se expresa como un número decimal, pero también se puede entender como un número entero y un residuo:
\begin{gather*}
  \text{dividendo} = \text{divisor} \times \text{cociente entero} + \text{residuo}
\end{gather*}
La división entera representa el cociente entero, y el módulo representa el residuo:
\begin{gather*}
  \text{dividendo} \operatorname{\texttt{//}} \text{divisor} = \text{cociente entero} \\
  \text{dividendo} \operatorname{\texttt{\%}} \text{divisor} = \text{residuo}
\end{gather*}

E.g., $11 / 4$ es igual a $2.75$, pero también se puede expresar como $11 = 4 \times 2 + 3$. En ese caso, la división entera de $11$ entre $4$ es $2$, que es simplemente tomar la parte entera del cociente\footnote{Nótese que esto no es lo mismo que aproximar.}, y el módulo de $11$ entre $4$ es $3$. 

Más intuitivamente, la división entera representa la cantidad de veces que el un número (el divisor) cabe completamente dentro de otro (el dividendo), y el módulo representa lo que sobra. Si se reparte una torta de 11 pedazos equitativamente entre 4 personas, a cada una le tocan 2 pedazos (el cociente entero) y sobran 3 pedazos que no le tocan a nadie (el módulo).

\subsubsection{Operadores de asignación compuesta}

Es común en programación modificar el valor de una variable operándola con algún valor. E.g.,
\begin{minted}{python}
  x = 10
  x = x + 5
\end{minted}
Que actualiza el valor de \mintinline{python}|x| a \mintinline{python}|15|. Existe una sintaxis más corta para escribir este tipo de operaciones, que es la siguiente:
\begin{minted}{python}
  x = 10
  x += 5
\end{minted}
El operador \mintinline{python}|+=| es un operador de asignación compuesta. Existe uno para cada operador aritmético y de cadenas de caracteres:
\begin{table}[h]
\centering
\begin{tabular}{|l|c|c|}
\hline
\textbf{Operador} & \textbf{Ejemplo} & \textbf{Asignación equivalente} \\
\hline
\mintinline{python}|+=| & \mintinline{python}|x += 5| & \mintinline{python}|x = x + 5| \\
\mintinline{python}|-=| & \mintinline{python}|x -= 3| & \mintinline{python}|x = x - 3| \\
\mintinline{python}|*=| & \mintinline{python}|x *= 2| & \mintinline{python}|x = x * 2| \\
\mintinline{python}|/=| & \mintinline{python}|x /= 4| & \mintinline{python}|x = x / 4| \\
\mintinline{python}|//=| & \mintinline{python}|x //= 3| & \mintinline{python}|x = x // 3| \\
\mintinline{python}|%=| & \mintinline{python}|x %= 2| & \mintinline{python}|x = x % 2| \\
\mintinline{python}|**=| & \mintinline{python}|x **= 3| & \mintinline{python}|x = x ** 3| \\
\mintinline{python}|+=| & \mintinline{python}|s += " Mundo"| & \mintinline{python}|s = s + " Mundo"| \\
\mintinline{python}|*=| & \mintinline{python}|s *= 3| & \mintinline{python}|s = s * 3| \\
\hline
\end{tabular}
\caption{Operadores de asignación con operador en Python}
\label{tab:operadores-asignacion-con-operador}
\end{table}


% TODO: Checkpoint. Implementar ambiente o barra linda de checkpoint.

A estas alturas, sabemos:
\begin{itemize}
  \item Qué es un valor
  \item Algunos tipos de dato primitivos en Python
  \item Qué es una variable y para qué sirve el operador de asignación.
  \item Qué es una expresión 
  \item Cuáles operaciones pueden aplicarse sobre los tipos de dato vistos.
  \item Cómo funcionan los operadores de asignación compuesta
\end{itemize}

Eso no permite hacer programas muy complejos, pero sí nos permite escribir programas que hagan cálculos y almacenen resultados. Ya en este punto, Python es más poderoso que la mayoría de calculadoras avanzadas. Probablemente, soporta cálculos más extensos que su calculadora.